% Generated draft: requires mathematical and source review.
\documentclass[11pt]{article}
\usepackage[margin=25mm]{geometry}
\usepackage{fontspec}
\setmainfont{Latin Modern Roman}
\usepackage{amsmath,amssymb,mathtools,unicode-math}
\setmathfont{Latin Modern Math}
\usepackage{enumitem,needspace,xurl,hyperref}
\setlength{\parindent}{0pt}
\setlength{\parskip}{5pt}
\setlength{\emergencystretch}{3em}
\newcommand{\N}{\mathbb{N}}
\newcommand{\Z}{\mathbb{Z}}
\newcommand{\Q}{\mathbb{Q}}
\newcommand{\R}{\mathbb{R}}
\newcommand{\C}{\mathbb{C}}
\newcommand{\F}{\mathbb{F}}
\newcommand{\A}{\mathbb{A}}
\newcommand{\PP}{\mathbb{P}}
\newcommand{\E}{\mathbb{E}}
\newcommand{\eps}{\varepsilon}
\newcommand{\dd}{\,\mathrm{d}}
\begin{document}
\section*{UnsolvedMath ID: 10400193}
\textbf{Source problem code:} catalogue entry
\subsection{Definitions and mathematical statement}
Retain \(q_\sigma\) as the source's notation for a family of quadratic forms. The supplied excerpt does not define \(\sigma\), the domains or value groups of these forms, their quadratic identities, or any nondegeneracy hypotheses.

\emph{Problem (F. Deloup).} Classify the monoid of isomorphism classes of the quadratic forms \(q_\sigma\), with operation induced by orthogonal sum:
\[
 \mathcal Q=\{[q_\sigma]\},\qquad
 [q_\sigma]+[q_\tau]=[q_\sigma\perp q_\tau].
\]
Here an isomorphism is an isomorphism of the underlying algebraic objects that preserves the quadratic form and any required structure. When the forms have additive domains \(A_\sigma,A_\tau\) and a common additive target, the usual orthogonal sum is
\[
 (q_\sigma\perp q_\tau)(x,y)=q_\sigma(x)+q_\tau(y),
 \qquad(x,y)\in A_\sigma\oplus A_\tau.
\]
This formula describes the operation conditionally; it does not identify the missing family. A classification should determine the isomorphism classes and their sum, for example through complete invariants or a presentation of the monoid, including its neutral element in the intended setup. The problem concerns actual isomorphism classes under orthogonal sum; it does not specify passing to a Witt group or a stable-equivalence quotient. Until \(q_\sigma\) is defined, a particular monoid cannot be reconstructed from this record.
\subsection{Short English statement}
Deloup asks to classify the quadratic forms denoted \(q_\sigma\), together with their orthogonal-sum operation. The family itself is not defined in the excerpt, so its underlying domains and quadratic-form conventions must be recovered before a concrete classification problem is fixed.
\subsection{Sources}
\begin{itemize}
\item Local UnsolvedMath registry, permanent ID 10400193, source code catalogue entry; attribution F. Deloup. The stored title is truncated after ``Problem 10.18 --- (F.''.
\item Ohtsuki, \emph{Problems on invariants of knots and 3-manifolds} (2002), Problem 10.18, PDF page 153. Supplied URL: \url{https://msp.org/gtm/2002/04/gtm-2002-04-024s.pdf}; recorded access 2026-07-29; extraction pdf-font-structure. The collection distinguishes 231 numbered bold problem, conjecture, and question headings.
\item The supplied status field is open and research\_classification is OPEN-TRIAGE, whereas the research summary begins ``Classification: PARTIAL-PROGRESS.'' No specific progress result is supplied. Current primary-literature status and redistribution terms for the publicly accessible author/publisher copy remain NEEDS\_REVIEW.
\end{itemize}
\end{document}
